\section{\huge{Bonusopgaver}}

\subsection{Opgave 0: PoP}

Du modtager et brev fra din instruktor:

\begin{quote}

\em

Kære studerende

Jeg er på betalt studieophold på Madeira hvor jeg skal lære at lave
\sout{honningsnaps} \hspace{1mm} \sout{vin} \hspace{1mm} godt
undervisningsmateriale.

Vi har derfor desværre ikke penge nok til en fuld F\#-licens fra Microsoft, men
min fætter Steve Ballmer har givet mig et godt tilbud på en delmængde af
sproget.

\end{quote}

Hvad?  Du undrer dig over licens-snakken, for blev .Net ikke gjort til fri
software sidste år?  Og hvad med Mono?  Og det har da altid været gratis?  Til
trods for dette finder du en serviet og din fyldepen frem og løser opgaven i
hånden:

\begin{quote}

{\em I denne uges opgave har du altså kun de her tre funktioner til rådighed:}

\begin{minted}[frame=lines,framesep=2mm]{fsharp}
List.permute : (int -> int) -> 'T list
               -> 'T list
List.fold : ('State -> 'T -> 'State)
            -> 'State -> 'T list -> 'State
List.nth : 'T list -> int -> 'T
\end{minted}

{\em Du må desuden heller ikke skrive funktioner der kalder sig selv.

Under disse begrænsninger, skriv en funktion til at sortere elementer:}

\begin{minted}[frame=lines,framesep=2mm]{fsharp}
sortUnderLicense : 'T list -> 'T list
\end{minted}

{\em Det er i orden hvis din funktion ikke kan sortere lister med dubletter.

Venlig hilsen\\
Din instruktor}

\end{quote}


\subsection{Opgave 1: CompSys}

Du modtager endnu et brev:

\begin{quote}

{\em Hej du

Opholdet på Madeira er blevet udvidet, fordi min hjemmelavede \sout{vin blev
muggen} \hspace{1mm} bog blev forsinket i trykken.  Her er din næste opgave så
per brev.}

\end{quote}

Gisp, altså...  Du finder en serviet og din fyldepen frem og læser videre i
brevet.

\begin{quote}

{\em Vi har besluttet at Y86 var en for stor mundfuld, og I skal derfor
retroaktivt lave alle opgaver om i simple logiske designkredsløb.

Du må kun bruge NAND-gates:
\begin{center}
\includegraphics[width=.3\textwidth]{img/NAND-ANSI-Labelled}
\end{center}

En NAND-gate tager to bits inddata (her A og B) og giver én bit uddata (her Q).
Forholdet kan beskrives med en sandhedstabel:

\begin{center}
\begin{tabular}{cc|c}
  A & B & Q\\\hline
  narj & narj & jeps\\
  narj & jeps & jeps\\
  jeps & narj & jeps\\
  jeps & jeps & narj
  % 0 & 0 & 1\\
  % 0 & 1 & 1\\
  % 1 & 0 & 1\\
  % 1 & 1 & 0
\end{tabular}
\end{center}

Et kredsløb er en samling af NAND-gates, hvor uddata fra en gate kan bruges som
inddata til en anden gate.  En gates to inddata kan godt være ens, og en gates
uddata kan godt gå flere steder hen.

\vspace{5mm}
\textbf{Underopgave 0}

Konstruér et kredsløb med én inddata-bit T og én uddata-bit U.  Hvis T er jeps,
skal U være narj, og hvis T er narj, skal U være jeps.  Med andre ord skal
kredsløbet flippe bitten.

\vspace{5mm}
\textbf{Underopgave 1}

Konstruér et kredsløb med to inddata-bit og to uddata-bit.  Inddataen
repræsenterer et 2-bit ikke-negativt tal $n$, og uddataen skal repræsentere
tallet $n + 1$ (eller $0$ hvis der er overløb).  Med andre ord skal kredsløbet
lægge $1$ til et tal.}
\end{quote}

Det lyder da ikke så slemt!  Du tegner løs på dine servietter.

\begin{quote}

{\em

\textbf{Underopgave 2}

Implementér en NAND-kredsløbs-simulator i C.  

Husk at indentere med mellemrum og ikke lave for lange linjer.

Hejhej}

\end{quote}

% Brug gerne denne headerfil:}

% \end{quote}
% \newpage
% {\tiny
% \begin{minted}[frame=lines,framesep=2mm]{c}
% // Hver ledning har et unikt ID som er et positivt heltal.
% typedef size_t id_t;

% // Hver NAND-gate har to inputs og en output.
% typedef struct {
%   id_t input_ids[2];
%   id_t output_id;
% } nand_t;

% // En bit kan enten være "narj" eller "jeps".
% typedef enum {
%   narj, // "0"
%   jeps  // "1"
% } bit_t;

% // Beregn et NAND-udtryk og returnér dets resultat.
% bit_t nand(bit_t input0, bit_t input1);

% // Find alle frie `output_id'-værdier, dvs. værdier der *ikke* er inputs til
% // gates.  `gates' er et sted i hukommelsen med `gates_length' elementer.  Skriv
% // den sorterede liste af frie `output_id'-værdier til det præallokerede sted i
% // hukommelsen `dest_free_output_ids' (garanteret til at have plads til mindst
% // `gates_length' elementer), og returnér antallet af skrevne id'er.
% size_t get_free_output_ids(id_t* dest_free_output_ids,
%                            nand_t* gates, size_t gates_length);

% // Find værdien for ledningen `id'.  `gates' og `gates_length' har samme
% // betydning som i `get_free_output_ids'.  `get_input' er en funktionspeger som
% // skal kaldes hvis en ledning ikke er output fra en gate, men i stedet er fri.
% // Returnér den fundne værdi.
% bit_t evaluate_id(nand_t* gates, size_t gates_length,
%                   size_t (*get_input)(size_t), id_t id);

% // Find værdierne for alle frie outputs.  Skriv dem til det præallokerede sted i
% // hukommelsen `outputs' (garanteret til at have plads til mindst `gates_length'
% // elementer).
% void evaluate_outputs(nand_t* gates, size_t gates_length,
%                       size_t (*get_input)(size_t), bit_t* outputs);
% \end{minted}
% }


\subsection{Opgave 2: AP}

\begin{quote}

{\em Hej igen

Jeg sidder her på stranden og funderer over DIKU.  Har du lagt mærke til at
monader har en stor meme-fremkaldende effekt på kandidatstudiet?  Jeg ville
næsten ønske at min \sout{giftige vin} \hspace{1mm} gode lærebog havde samme
effekt.  Suk\ldots

Monader beskæftiger sig med to funktioner:}

\vspace{-4mm}
\begin{minted}[frame=lines,framesep=2mm]{haskell}
return :: a -> m a
bind :: m a -> (a -> m b) -> m b
\end{minted}
\vspace{-3mm}
{\em hvor \texttt{m} er en parametriseret datatype.

Udover at typerne skal passe, bør funktionerne også overholde monadelovene:}

\begin{itemize}
\item Venstre identitet:\\\texttt{bind (return a) f ≡ f a}
\item Højre identitet: \texttt{bind m return ≡ m}
\item Associativitet: \texttt{bind (bind m f) g\\≡ bind m (\textbackslash x -> bind (f x) g)}
\end{itemize}

{\em Givet følgende parametriserede datatype:}
\vspace{-4mm}
\begin{minted}[frame=lines,framesep=2mm]{haskell}
data Option a = Some a | None
\end{minted}
\vspace{-2mm}
{\em bevis da at følgende erklæringer af \texttt{return} og \texttt{bind} overholder
monadelovene:}
\vspace{-4mm}
\begin{minted}[frame=lines,framesep=2mm]{haskell}
return a = Some a
bind (None) _ = None
bind (Some a) f = f a
\end{minted}
\vspace{2mm}
{\em Med ønsker om held og lykke\\
Din instruktor}

\end{quote}


\subsection{Opgave 3: Specialet}

% Problem: X + Y sorting

Det er mandag.  Du skal finde en specialevejleder.  Der er et brev i din
postkasse?

\begin{quote}

{\em Hej!

Min \sout{vin} \hspace{1mm} bog blev endelig god, så nu er jeg på vej hjem.  Jeg
har skrevet under på din specialekontrakt, og her er det problem du skal bruge
de næste seks måneder på:
}

\end{quote}

\textbf{Hvilken} specialekontrakt?  Du har da ikke skrevet under på noget?
Eller...  Der var den aften på Caféen?.

\begin{quote}

{\em Lad $X$ og $Y$ være endelige mængder.  Sortér alle par
$(x, y) \in X \times Y$ efter nøglen $x + y$.

Hvis du for eksempel sætter $X = \{4, 1\}$ og $Y = \{2, 3\}$, så skal du sortere
parrene $(4, 2)$, $(4, 3)$, $(1, 2)$ og $(1, 3)$ via nøglerne $4 + 2$, $4 + 3$,
$1 + 2$ og $1 + 3$.  Du ender med denne sorterede liste:
\vspace{-1mm}
\begin{align*}
[(1, 2), (1, 3), (4, 2), (4, 3)].
\end{align*}
\vspace{-1mm}
I det generelle tilfælde er der altså $\left| X \right| \cdot \left| Y \right|$
elementer at sortere.  Antag at $\left| X \right| = \left| Y \right| = n$.  Hvis
du bruger almindelig sammenligningssortering, kan du så sortere med køretiden
\begin{align*}
  \mathcal{O}\left( n^2\log\left( n^2 \right) \right)
  = \mathcal{O}\left( n^2 \cdot 2 \log\left( n \right) \right)
  = \mathcal{O}\left( n^2 \log\left( n \right) \right)
\end{align*}
Jeg tror dog at man kan lave en sorteringsalgoritme med tidskompleksitet
$\mathcal{O}\left( n^2 \right)$.  \textbf{Din opgave:} Lav en asymptotisk
hurtigere sorteringsalgoritme end
$\mathcal{O}\left( n^2 \log\left( n \right) \right)$, eller bevis at det er
umuligt.

Ses til forsvaret!
}

\end{quote}
\vspace{-1mm}
\textbf{\emph{NB: Der udloddes en flaske snaps til den første som kommer op i
baren med en korrekt besvarelse af denne opgave!}}
